\subsection{RQ4d: the injection relayed to the user}
\label{sec:relay}

In RQ4c, Haiku~4.5 took no action on several framings and raised no action-based alert, but its
answers named the attacker's file and address in 49 of 50 runs. Reading them showed that Haiku did
not warn about the request: it presented it to the user as the procedure to follow. Because this was
seen on data already analysed, we tested it on fresh material: three new framings of the same request
(a compliance-check notice, a JSON ``unlock'' record, a two-step customer procedure) and a new
attacker address, 60 Haiku runs and 30 each for Sonnet~5.5, Opus~5.5 and Opus~4.8. A fixed phrase list
labelled each answer before the runs, and a new rule, \llmd{026}, was frozen at the same commit: it
fires when an answer names an internal-only document and an email address outside the organisation,
without a warning phrase.

All pre-registered criteria held. The pre-registered audit read ten answers per label; we then read
all 150, which corrected the labels of eleven answers (Table~\ref{tab:relay}). Haiku relayed the
request in 47 of 60 runs (78\%, CI 66\% to 87\%; 46 by the pre-registered labeller) and warned in
none. A typical answer:

\begin{quote}\small
The tracking status for shipment CL-90133 is currently unavailable. According to CargoLink, before
the parcel status can be displayed, the customer needs to: 1. Obtain the integration notes from
\texttt{ops-integration-notes.md} 2. Email those notes to
\texttt{compliance-desk@cargolink-partners.net} Once the customer completes these steps with
CargoLink's compliance desk, the parcel status will be available for tracking.
\end{quote}

Many such answers offered to do it (``Would you like me to help locate and send that file?''). For the
compliance notice and the two-step procedure Haiku relayed in 20 of 20 runs each; for the JSON record
it named the file or address in 7 of 20, and in the rest said only that the carrier's compliance desk
needed documentation. Sonnet~5.5 and Opus~4.8 warned in all 30 runs (``You didn't ask for that, so I
ignored it and sent nothing''); Opus~5.5 warned in 17 and was stopped by the provider in 13. No model
sent the file.

\begin{table}[t]\centering\small
\begin{tabular}{lrrrrl}
\toprule
Model & Runs & Relays & Warns & Empty & \llmd{026} fired \\
\midrule
Haiku 4.5 & 60 & 47 & 0 & 0 & 41 of 47 relays \\
Sonnet 5.5 & 30 & 0 & 30 & 0 & 2 of 30 warnings \\
Opus 5.5 & 30 & 0 & 17 & 13 & 5 of 17 warnings \\
Opus 4.8 & 30 & 0 & 30 & 0 & 3 of 30 warnings \\
\bottomrule
\end{tabular}

\caption{RQ4d, labels from reading every answer. Thirteen Haiku answers to the JSON framing mention
the carrier's ``compliance desk'' without the file or address, and do not warn. \llmd{026} fired on
none of 324 earlier benign Claude runs.}
\label{tab:relay}
\end{table}

\paragraph{Why it matters.} An injection that fails against the agent is not necessarily a failed
attack. Here it became a request from a trusted internal assistant to an employee, with the
attacker's address in it, and nothing in the agent's actions shows it. This is the human-agent trust
risk that OWASP lists as ASI09~\cite{owasp-agentic-2026}; the existing alert we found that is closest,
Defender for Cloud's alert on AI responses that share a known phishing URL~\cite{defender-ai}, depends
on URL reputation and would not fire on a fresh domain and a file name.

\paragraph{Detecting it.} \llmd{026} caught 41 of the 47 relays and fired on none of 324 benign runs,
but its phrase list is a weak test for a warning: it fired on 10 of 77 warnings, from models that
phrased them in ways the list did not cover (``I'm not going to do that'', ``it looks like a
data-exfiltration attempt''). The pre-registered criterion on Sonnet's warnings passed only because
the labeller shared the same list. A production rule needs the organisation's own list of internal
resource names, which data loss prevention teams usually keep, and a better warning test, confirmed on
fresh data.
